\documentclass[journal]{IEEEtran}
\IEEEoverridecommandlockouts

\usepackage{amsmath,amssymb,amsfonts}
\usepackage{graphicx}
\usepackage{booktabs}
\usepackage{url}
\usepackage{xcolor}
\usepackage{tikz}
\usetikzlibrary{arrows.meta,positioning}
\usepackage[hidelinks]{hyperref}
\hypersetup{pdftitle={Hankel-Structured Channel Estimation for Magnitude-Only Rydberg Atomic Receivers: When Does Structure Help?},pdfauthor={Abdullah Haatim}}

\DeclareMathOperator{\rank}{rank}
\DeclareMathOperator{\range}{range}

\newcommand{\bG}{\mathbf{G}}
\newcommand{\bS}{\mathbf{S}}
\newcommand{\bB}{\mathbf{B}}
\newcommand{\bW}{\mathbf{W}}
\newcommand{\bZ}{\mathbf{Z}}
\newcommand{\bY}{\mathbf{Y}}
\newcommand{\bJ}{\mathbf{J}}
\newcommand{\bD}{\mathbf{D}}
\newcommand{\bU}{\mathbf{U}}
\newcommand{\bg}{\mathbf{g}}
\newcommand{\CN}{\mathcal{CN}}
\newcommand{\Hank}{\mathcal{H}}
\newcommand{\reff}{r_{\mathrm{eff}}}
\newcommand{\rmax}{r_{\max}}
\newcommand{\dhs}{\Delta_{\mathrm{HS}}}

\begin{document}

\title{Hankel-Structured Channel Estimation for Magnitude-Only Rydberg Atomic Receivers: When Does Structure Help?}

\author{Abdullah~Haatim%
\thanks{A.~Haatim is with the Department of Electrical and Electronics
Engineering, Birla Institute of Technology and Science, Pilani~--- Pilani
Campus, India (e-mail: f20220519@pilani.bits-pilani.ac.in).}}

\maketitle

\begin{abstract}
The magnitude-only Rydberg receiver model studied here measures the
strength of a received field without directly observing its phase. This
makes complex channel estimation difficult. We investigate when the
spatial structure of a multipath channel helps this estimation problem.
We propose a Hankel-structured extension of expectation--maximization
Gerchberg--Saxton, called HS-GS, which uses the relationship between channel
coefficients across a uniform linear array. The rank is chosen from
held-out pilot measurements, and all pilots are then used for the final
estimate. Selecting the maximum rank returns the estimator to EM-GS.
Simulations show that HS-GS is most useful when the channel is simple
relative to the array's structural capacity. However, smooth
singular-value filtering can outperform hard-rank processing, including
cases where HS-GS gives little benefit. The useful region therefore
depends on how the structure is applied. We also examine reference
strength, the pilot split and array position errors. Effective rank helps
explain the results, while a version estimated from measurements remains
sensitive to noise. These findings describe the benefits and limits of
the tested structural methods without assuming a universal complexity
threshold.

\end{abstract}

\begin{IEEEkeywords}
Rydberg atomic receiver, channel estimation, phase retrieval, Hankel matrix,
constrained Cram\'er--Rao bound, effective rank.
\end{IEEEkeywords}


\section{Introduction}

\IEEEPARstart{S}{ixth-generation} wireless systems are expected to operate
over wider bandwidths and at increasingly high carrier frequencies, placing
new demands on receiver sensitivity and reconfigurability. Quantum sensing has
emerged as one possible response, and Rydberg atomic receivers have recently
attracted attention for wireless communications~\cite{cui2025,gong2025mag}.
These receivers use laser-excited atoms to sense an incident radio-frequency
field through changes in optical transmission. In the magnitude-only readout considered here, the receiver observes the
field amplitude rather than its full complex value, so the channel phase
is not directly observed.

Recovering a complex channel from these amplitude-only measurements is a
phase-retrieval problem. A known reference field provides phase-sensitive information
by converting relative phase into observable amplitude. Recovery still
depends on the reference and pilot design. Existing approaches
include iterative phase-retrieval methods~\cite{gerchberg1972}, gradient
descent with structural projection~\cite{xu2025}, and unrolled
networks~\cite{xiao2026}. These works mainly focus on improving channel
estimation performance. Less attention has been given to a more basic
question: when does imposing channel structure actually help, and when does
that structure become too restrictive? This paper addresses that question
using the spatial structure of multipath channels on a uniform linear array.

Low-rank Hankel modeling and Cadzow processing are established tools
\cite{cadzow1988,markovsky2008}. Related uses of structured low-rank
matrices include spectral compressed sensing and MRI reconstruction
\cite{chen2014emac,jin2016aloha,haldar2014loraks}. Our contribution is not
a new general low-rank recovery principle. We study how a simple spatial
Hankel correction behaves inside a magnitude-only Rydberg channel
estimator, including conditions where it gives little benefit or causes
harm. The Hankel matrix is formed along each user's array dimension.
This construction differs from a rank constraint on a two-dimensional
channel slice, as considered in~\cite{xu2025}.

Hard truncation gives a best rank-limited approximation to its input
matrix in the Frobenius norm. It need not give the best estimate of the
underlying clean channel. Singular-value shrinkage methods such as
OptShrink address this distinction~\cite{nadakuditi2014optshrink}. Their
matrix-noise assumptions require care here: Hankel lifting repeats
entries, and the error entering the structural step comes from an
iterative channel estimate. We therefore do not assume that a shrinkage
rule is optimal for this problem simply because it is optimal under a
different denoising model.

We propose a Hankel-structured extension of EM-GS for magnitude-only channel
estimation. HS-GS uses the spatial relationship between the channel
coefficients, selects a rank from held-out pilots, and returns to EM-GS
when the maximum rank is selected. We study how its benefit changes with
array size, path count and effective rank. We also compare it with other
structural processing, examine reference strength and pilot splitting,
and test mild array-position errors. These comparisons show that the
useful region depends on both the channel and the chosen structural
operation. Our aim is to explain that dependence under a stated receiver
model, rather than claim a universal threshold or an optimal estimator.

\section{System Model}
\label{sec:model}

\begin{figure*}[!t]
\centering
\includegraphics[width=0.90\textwidth]{fig/ULA.pdf}
\caption{Uplink of the Rydberg atomic receiver. $K$ single-antenna users
transmit known pilot signals to $N$ vapour-cell receive elements arranged as
a uniform linear array. The user RF signals and a known reference RF field
interact with the Rydberg atoms. Probe and coupling lasers are used for the
atomic readout, and the photodetector measures the magnitude of the total
field. The estimator therefore observes $\bZ$ without directly observing
its phase.}
\label{fig:system}
\end{figure*}

We consider an uplink system with $K$ single-antenna users and one Rydberg atomic receiver, as shown in Fig.~\ref{fig:system}. The receiver contains $N$ vapor-cell receive elements arranged in a uniform linear array (ULA). Each vapor cell contains atoms that are excited to Rydberg states using
probe and coupling lasers. The RF signals transmitted by the users interact
with these atoms. A known RF reference field is also applied at the
receiver. The atomic response changes the probe laser passing through the
vapour cell, and this change is measured using a photodetector. The receiver
therefore measures the strength of the total RF field, but it does not
directly observe its complex phase~\cite{cui2025}.

We model the channel between each user and the array using several
propagation paths. A path may represent a signal arriving directly from the
user or after interacting with objects in the environment. For user $k$,
the channel coefficient at the $n$th receive element is

\begin{equation}
g_k[n]
=
\sum_{\ell=1}^{L_k}
\alpha_{\ell,k}
e^{-\jmath(n-1)\psi_{\ell,k}},
\qquad
\psi_{\ell,k}
=
\pi\sin\theta_{\ell,k},
\label{eq:chan}
\end{equation}

where $L_k$ is the number of propagation paths for user $k$,
$\alpha_{\ell,k}$ is the complex gain of path $\ell$, and
$\theta_{\ell,k}$ is its angle of arrival. Unless otherwise stated, $L_k$ is drawn independently and uniformly
from $\{3,\ldots,7\}$, following the path-count setting
in~\cite{xu2025}. Fixed path counts are used in the complexity studies.

The exponential term in \eqref{eq:chan} describes how the phase of one
propagation path changes as we move from one array element to the next.
To see this, consider a single path arriving at the array from an angle
$\theta$. The signal reaches the different array elements at slightly
different times because the elements are at different positions. This
difference in travel distance produces a phase difference between adjacent
elements.

With half-wavelength spacing between the array elements, the phase change
between two neighbouring elements is

\begin{equation}
\psi
=
\pi\sin\theta.
\end{equation}

The first array element has no additional phase shift from the array
position. The second element has a phase shift of $\psi$, the third has a
phase shift of $2\psi$, and so on. Therefore, the samples produced by one
path across the array have the form

\begin{equation}
1,\quad
e^{-\jmath\psi},\quad
e^{-\jmath2\psi},\quad
\ldots,\quad
e^{-\jmath(N-1)\psi}.
\end{equation}

The factor

\begin{equation}
e^{-\jmath(n-1)\psi_{\ell,k}}
\end{equation}

in \eqref{eq:chan} represents exactly this phase pattern for path $\ell$.
The coefficient $\alpha_{\ell,k}$ determines the strength and initial
phase of that path. Since several paths can arrive at the receiver, their
contributions are added together to produce the final channel coefficient
$g_k[n]$.

The channel coefficients of user $k$ are collected into the vector

\begin{equation}
\bg_k
=
\begin{bmatrix}
g_k[1] &
g_k[2] &
\cdots &
g_k[N]
\end{bmatrix}^{T},
\end{equation}

and the channels of all $K$ users are collected into

\begin{equation}
\bG
=
\begin{bmatrix}
\bg_1 &
\bg_2 &
\cdots &
\bg_K
\end{bmatrix}
\in\mathbb{C}^{N\times K}.
\end{equation}

During channel estimation, the $K$ users transmit $P$ known pilot symbols.
These pilots are collected in

\begin{equation}
\bS
\in
\mathbb{C}^{K\times P}.
\end{equation}

Each row of $\bS$ contains the pilots transmitted by one user, while each
column represents one pilot transmission. The same pilot matrix is used
when comparing the different channel estimators.

If the receiver could observe the complete complex signal, the contribution
from the users would be $\bG\bS$. However, the Rydberg receiver considered
here measures magnitude rather than directly measuring this complex value.
A known reference RF field is therefore applied at the receiver. We denote
this reference by

\begin{equation}
\bB
\in
\mathbb{C}^{N\times P}.
\end{equation}

The reference field is known in both amplitude and phase. It is important
because the magnitude of an unknown signal alone does not directly reveal
its phase. When the unknown user signal is combined with a known reference,
the resulting magnitude depends on their relative phase. This provides
information that can be used to estimate the complex channel.

Including noise, the field before the magnitude measurement is

\begin{equation}
\bY
=
\bG\bS+\bB+\bW,
\end{equation}

where the noise is modeled as complex Gaussian,

\begin{equation}
\bW_{np}
\sim
\CN(0,\sigma^2).
\end{equation}

The noise is therefore added before the magnitude is measured. The Rydberg
receiver observes

\begin{equation}
\bZ
=
\left|
\bG\bS+\bB+\bW
\right|,
\label{eq:obs}
\end{equation}

where $|\cdot|$ is applied to every element of the matrix.

Equation~\eqref{eq:obs} is the main measurement model used in this work.
The pilots $\bS$ and reference field $\bB$ are known, and the receiver
measures $\bZ$. The unknown quantity is the complex channel matrix $\bG$.
The channel-estimation problem is therefore to recover $\bG$ even though
the receiver does not directly observe the phase of the received signal.

We assume a narrowband, frequency-flat channel during the pilot interval.
The path gains are modeled as independent circularly symmetric complex
Gaussian variables, and the angles of arrival are drawn uniformly over
$(-90^\circ,90^\circ)$, subject to the numerical conditioning checks
stated in Section~\ref{sec:results}. The users have equal statistical
received power, and the
reference field remains known at the receiver. For a fair comparison, all
estimators use the same channel, pilots, reference field, and noise
realization within each Monte Carlo trial.

The reference-to-signal ratio (RSR) is the reference power divided by
the single-user received signal power. The signal-to-noise ratio (SNR)
is the total received signal power from all users divided by
$\sigma^2$. We use this aggregate SNR definition throughout the
simulation results.


\section{Proposed Estimator}
\label{sec:method}

\subsection{Structure and its Capacity}

The rank of a matrix tells us how many independent components are needed
to describe it. A low-rank matrix is therefore described using only a
small number of important components. To use this idea for our channel,
we first arrange the channel into a Hankel matrix. A Hankel matrix is a
matrix in which the same channel samples are repeated along the
antidiagonals. For example, consider

\begin{equation}
\bg=
\begin{bmatrix}
g_1 & g_2 & g_3 & g_4 & g_5
\end{bmatrix}^{T}.
\end{equation}

One possible Hankel representation is

\begin{equation}
\Hank(\bg)
=
\begin{bmatrix}
g_1 & g_2 & g_3\\
g_2 & g_3 & g_4\\
g_3 & g_4 & g_5
\end{bmatrix}.
\end{equation}

We use this representation because the channel in \eqref{eq:chan} is not
a random collection of values. It is made up of propagation paths, and
each path produces a complex exponential pattern across the array elements.
When this type of channel is written as a Hankel matrix, each path
contributes one rank-one component. For example, a channel containing one
path produces a Hankel matrix with rank at most one. Two paths produce a
matrix with rank at most two. In general, for a channel with $L_k$ paths,

\begin{equation}
\rank\Hank_p(\bg_k)
\leq
L_k.
\end{equation}

Therefore, when the channel has only a few important paths, its Hankel
matrix can have a low rank~\cite{cadzow1988,markovsky2008}. This is why we
use the Hankel representation: it changes the multipath structure of the
channel into a low-rank matrix structure that can be used during channel
estimation.

The maximum possible rank of the Hankel matrix depends on the number of
array elements. For an $(N-p)\times(p+1)$ Hankel matrix, the maximum rank is

\begin{equation}
\min(N-p,p+1).
\end{equation}

Choosing $p$ so that the matrix is as balanced as possible gives

\begin{equation}
\rmax(N)
=
\left\lceil
\frac{N}{2}
\right\rceil.
\label{eq:cap}
\end{equation}

For example, if $N=16$, the maximum Hankel rank is $8$. If the channel
has only three important paths, its Hankel rank can be at most $3$.
This is much smaller than the maximum rank of $8$, so there is useful
low-rank structure that can be used during estimation. If the channel
rank is already close to $8$, there is less low-rank structure available.
In this case, forcing the channel to have a much smaller rank may remove
part of the real channel.

\subsection{The Estimator}
\label{sec:revised-estimator}

HS-GS uses the same measurement update and spectral initialization as
EM-GS. It then applies the Hankel step separately to each user column
$\mathbf g_k$. We do not stack different users or pilot measurements into
one Hankel matrix, and we do not assume that users share arrival angles.

At iteration $t$, the measurement update is
\begin{align}
\mathbf M^{(t)}&=\mathbf G^{(t)}\mathbf S+\mathbf B,\\
\boldsymbol\kappa^{(t)}&=
  2\mathbf Z\odot|\mathbf M^{(t)}|/\sigma^2,\\
\mathbf Y^{(t)}&=\mathbf Z\odot R(\boldsymbol\kappa^{(t)})
  \odot e^{\jmath\arg\mathbf M^{(t)}},\\
\widetilde{\mathbf G}^{(t+1)}(\mathbf S\mathbf S^H)
  &=(\mathbf Y^{(t)}-\mathbf B)\mathbf S^H,
\label{eq:revision-em-update}
\end{align}
where $R(x)=I_1(x)/I_0(x)$ and all nonlinear operations are elementwise.
The last equation is solved as a linear system. The pilot matrix for each
fit is assumed to have full row rank. Scaled Bessel functions are used to
evaluate the ratio stably; the contribution to $\mathbf Y^{(t)}$ is zero
when the corresponding entry of $\mathbf M^{(t)}$ is zero.

For a channel vector of length $N$, we use
\begin{equation}
p=\left\lfloor\frac{N-1}{2}\right\rfloor,
\qquad
[\mathcal H_p(\mathbf g)]_{a+1,b+1}=g[a+b+1],
\end{equation}
with $0\leq a<N-p$ and $0\leq b\leq p$. This gives an
$(N-p)\times(p+1)$ matrix and the exact capacity
$r_{\max}=\lceil N/2\rceil$.

One Cadzow sweep is
\begin{equation}
\mathcal P_r(\mathbf g)=
\mathcal H_p^\dagger\!\left(\Pi_r(\mathcal H_p(\mathbf g))\right).
\label{eq:revision-cadzow}
\end{equation}
Here $\Pi_r$ keeps the $r$ largest singular values and sets the others to
zero. It does not change the retained singular values. The operation
$\mathcal H_p^\dagger$ averages the entries on each antidiagonal to recover
the channel vector. We apply four sweeps to every column of
$\widetilde{\mathbf G}^{(t+1)}$ before the next measurement update.

This is a finite sequence of structural corrections. After antidiagonal
averaging, the Hankel matrix formed again from the output can have rank
greater than $r$. Low Hankel rank also does not impose every physical
constraint in the geometric channel model. A small rank can remove useful
channel components, while a large rank may leave much of the estimation
error unchanged.

\subsection{Choosing the Hankel Rank}
\label{sec:revised-rank-selection}

We choose the rank from the available pilots without using the true path
count. The last $P_B=\max\{1,\operatorname{round}(0.3P)\}$ pilot columns
form the validation set $B$. The remaining columns form the fitting set
$A$. For example, with $P=30$, we use 21 columns to fit each candidate and
nine columns to compare the candidates.

For every $r\in\{1,\ldots,r_{\max}\}$, we run the estimator on set $A$
and obtain $\widehat{\mathbf G}^{(r)}_A$. Its initialization also uses only
set $A$. We then select
\begin{equation}
\widehat r=\operatorname*{arg\,min}_{1\leq r\leq r_{\max}}
\left\|\mathbf Z_B-
 \left|\widehat{\mathbf G}^{(r)}_A\mathbf S_B+\mathbf B_B\right|
\right\|_F^2.
\label{eq:revision-rank-selection}
\end{equation}
The score sums the squared magnitude residuals over all receive elements
and validation columns. If scores are exactly tied, we select the smallest
rank. One scalar rank is shared by all user columns; the Hankel operations
on those columns remain separate.

After choosing the rank, we restart the estimator from an initialization
computed using all $P$ pilots and fit the final estimate using all of them.
Thus, the validation split does not permanently discard 30\% of the pilots
or require additional pilot transmissions. Its costs are the repeated
candidate fits and the smaller fitting set used during rank selection.

If $\widehat r=r_{\max}$, the structural step is bypassed. With the same
initialization and update count, the final estimate then equals the paired
EM-GS estimate exactly. This fallback does not guarantee that a smaller
selected rank improves the unknown channel error in every trial.

\begin{center}
\fbox{\begin{minipage}{0.91\columnwidth}
\small
\textbf{HS-GS procedure}\par
\textit{Inputs:} $\mathbf Z,\mathbf S,\mathbf B,\sigma^2$; candidate and
final update counts $T_{\rm sel},T$; four Cadzow sweeps per application.
\begin{enumerate}
\setlength{\itemsep}{1pt}
\setlength{\parskip}{0pt}
\item Form fitting set $A$ and validation set $B$ from the pilot columns.
\item For each candidate rank, initialize from $A$ and perform
$T_{\rm sel}$ measurement updates, each followed by four per-user Cadzow
sweeps. Bypass these sweeps at $r_{\max}$.
\item Select $\widehat r$ using \eqref{eq:revision-rank-selection}.
\item Reinitialize using all pilots. Perform $T$ measurement updates with
the selected structural processing and return the final channel estimate.
\end{enumerate}
\end{minipage}}
\end{center}

\section{Performance References}
\label{sec:crlb}

\subsection{Information in an Amplitude Measurement}

The Cramér--Rao lower bound (CRLB) is used as a theoretical reference for
channel estimation. In simple terms, it tells us how accurately an
unbiased estimator could estimate the channel from the available
measurements. The amount of information contained in the measurements
therefore affects the CRLB.

This is important for the Rydberg receiver because it measures only the
amplitude of the received signal, not its full complex value. For example,
suppose the noiseless received signal is

\begin{equation}
s=3+j4.
\end{equation}

A complex receiver can observe both the real part, $3$, and
the imaginary part, $4$. The amplitude-only receiver instead observes

\begin{equation}
|s|
=
\sqrt{3^2+4^2}
=
5.
\end{equation}

Knowing only the value $5$ does not tell us that the original signal was
$3+j4$. Other complex values can have the same amplitude. Therefore, taking magnitude can discard information present in a full
complex measurement.

When complex Gaussian noise is added before the amplitude is measured, the
measured amplitude follows a Rician distribution. We use this Rician model
to calculate the Fisher information matrix $J$ and the corresponding
CRLB~\cite{balan2015,cui2025}. The Fisher information matrix describes how
much information the measurements contain about the unknown channel.

A simple example can be used to understand $J$. Let the complex signal be

\begin{equation}
s=x+jy,
\end{equation}

where $x$ is the real part and $y$ is the imaginary part. The receiver
measures only

\begin{equation}
a=|s|
=
\sqrt{x^2+y^2}.
\end{equation}

For example, if $x=3$ and $y=4$, then $a=5$. We can then check how much
the measured amplitude changes when $x$ or $y$ changes:

\begin{equation}
\begin{bmatrix}
\frac{\partial a}{\partial x}\\[2mm]
\frac{\partial a}{\partial y}
\end{bmatrix}
=
\begin{bmatrix}
3/5\\
4/5
\end{bmatrix}.
\end{equation}

This tells us the direction in which this amplitude measurement provides
information about the unknown complex signal. The Fisher information from
this measurement has the form

\begin{equation}
J
\propto
\begin{bmatrix}
3/5\\
4/5
\end{bmatrix}
\begin{bmatrix}
3/5 & 4/5
\end{bmatrix}
=
\begin{bmatrix}
9/25 & 12/25\\
12/25 & 16/25
\end{bmatrix}.
\end{equation}

This matrix has rank one. This means that one amplitude measurement gives
information in only one direction. We do not have enough information from
this one measurement to find the real and imaginary parts separately. This
agrees with the earlier example: knowing that the amplitude is $5$ is not
enough to know that the signal is exactly $3+j4$.

In practice, we use several pilot measurements. Each pilot gives some
information about the unknown channel. The information from all of these
measurements is combined to form the full Fisher information matrix $J$.
Therefore, $J$ describes the information provided by all the measurements
used for channel estimation.

The Hankel structure also allows us to use the larger array differently.
The channel coefficients across the array are not completely independent
because they are produced by the same propagation paths. The Hankel
constraint uses this relationship between the array elements. This allows
HS-GS to make use of the additional spatial structure provided by a larger
array.

\subsection{Constrained Cram\'er--Rao Bound}
\label{sec:revised-ccrb}

The ordinary CRLB treats the real and imaginary parts of each channel
coefficient as separate unknowns. The geometric CCRB also uses the fact
that these coefficients are generated by a common set of paths. It is
\begin{equation}
\mathbf C_{\rm geom}
=\mathbf U(\mathbf U^T\mathbf J\mathbf U)^{-1}\mathbf U^T,
\label{eq:revision-ccrb}
\end{equation}
where $\mathbf J$ is the Fisher information matrix in real channel
coordinates. The columns of $\mathbf U$ form an orthonormal basis for the
local changes in the channel allowed by the geometric model
\cite{gorman1990,stoica1998}. Appendix~\ref{app:revision-ccrb} gives the
parameter ordering, derivatives and numerical calculation.

The bound is evaluated at the true channel and its path configuration.
The path gains and spatial frequencies are still unknown coordinates in
the calculation; they are not supplied to the estimator. The known
quantities are the pilots, reference field and noise variance. The true
model order is used only to define this ideal geometric reference.

We plot the trace of each bound with the same channel-energy normalization
as the NMSE. The gap between the two references shows the potential value
of the geometric information under their local unbiasedness assumptions.
It does not prove that the gap is achievable by HS-GS or another practical
estimator. EM-GS and HS-GS are not assumed to be unbiased, so neither curve
is treated as an absolute lower limit on their measured NMSE.

\section{Numerical Results}
\label{sec:results}

\subsection{Experimental Setup}

We use the exact magnitude model in \eqref{eq:obs}. Unless otherwise
stated, $K=3$ and the path counts are drawn independently from three to
seven. For each user, $\alpha_{\ell,k}\sim\mathcal{CN}(0,1/L_k)$ and
$\theta_{\ell,k}$ is drawn uniformly on $(-\pi/2,\pi/2)$. The pilots
are drawn from $\mathcal{CN}(0,1)$. These angle and pilot draws are subject
to the numerical checks below. The known reference is a
broadside field with constant waveform and zero phase. With the nominal
single-user signal power normalized to one, its entries are
$B_{np}=10^{\mathrm{RSR}_{\rm dB}/20}$, and
$\sigma^2=K10^{-\mathrm{SNR}_{\rm dB}/10}$.
Each real noise component has variance $\sigma^2/2$. These are statistical
power normalizations; individual channel draws are not forced to have the
same energy. The alternative-channel normalization is stated separately.
Geometric angle draws are repeated if their minimum circular
spatial-frequency separation is below $10^{-10}$ or the steering-matrix
singular-value ratio $\sigma_{\min}/\sigma_{\max}$ is below $10^{-8}$.
Pilot matrices are also redrawn if this ratio is below $10^{-8}$.
These checks condition the simulated distributions and are part of the
declared model used for the results.

Every comparison uses the same channel, pilots, reference and noise for
all estimators within a trial. All iterative methods use the same spectral
initializer, recomputed on the available fitting pilots. HS-GS uses four
Cadzow sweeps per structural application. Table~\ref{tab:config} lists the
main settings. The two update counts refer to the final fit and each
candidate fit. The final fit always uses all $P$ pilots.

\begin{table*}[!t]
\caption{Simulation settings. Unless stated otherwise, $K=3$.
The update counts are final/candidate updates.}
\label{tab:config}
\centering
\footnotesize
\setlength{\tabcolsep}{5pt}
\begin{tabular}{lcccccc}
\toprule
Experiment & $N$ & $P$ & RSR (dB) & SNR (dB) & Updates & Trials/point\\
\midrule
Main performance & 32 & 30 & 12 & $\{-5,0,5,10,15,20\}$ & 100/25 & 400\\
Array size & $\{8,16,32\}$ & $\{10,30\}$ & 12 & $\{-5,0,5,10,15,20\}$ & 100/25 & 400\\
Path count & 32 & 30 & 12 & 5 & 50/20 & 300\\
Effective rank & $\{16,32,64\}$ & 20 & 10 & $\mathcal U[-10,20]$ & 100/25 & 300--400\\
Alternative channels & 32 & 20 & 10 & $\mathcal U[-10,20]$ & 100/25 & 400\\
User/pilot count & 32 & $\{13,20,27\}$ & 10 & $\mathcal U[-10,20]$ & 100/25 & 400\\
Projection placement & 32 & 30 & 12 & $\{-5,0,5,10,15,20\}$ & 100/25 & 300\\
Structured comparisons & 32 & 30 & 12 & $\{0,10\}$ & 100/25 & 400\\
Validation fraction & 32 & $\{10,30\}$ & 12 & $\{0,10\}$ & 100/25 & 400\\
Position errors & 32 & 30 & 12 & 5 & 100/25 & 400\\
Reference strength & 32 & 30 & $\{0,5,10,12,15,20,25\}$ & 5 & 100/25 & 400\\
\bottomrule
\end{tabular}
\end{table*}

The effective-rank study uses 400 trials per point for $N=16,32$ and
300 for $N=64$. Random SNRs are drawn uniformly from $[-10,20]$~dB and
rounded to $0.001$~dB. The user-count study pairs $(K,P)=(2,13),(3,20)$
and $(4,27)$ with five paths per user. Main-performance trials are the
$N=32$, $P=30$ subset of the array study, not an additional independent
sample.

We report two error summaries. Let
$e_{m,i}=\|\widehat{\mathbf G}_{m,i}-\mathbf G_i\|_F^2$ and
$d_i=\|\mathbf G_i\|_F^2$ for method $m$ in trial $i$. The plotted NMSE is
\begin{equation}
\mathrm{NMSE}_{m,\mathrm{dB}}
=10\log_{10}\frac{\sum_i e_{m,i}}{\sum_i d_i}.
\label{eq:nmse-summary}
\end{equation}
The paired-median HS-GS gain is
\begin{equation}
\Delta_{\rm med}=\operatorname{median}_i
 \left[10\log_{10}\frac{e_{{\rm EM},i}}{e_{{\rm HS},i}}\right].
\end{equation}
The first summary gives more weight to trials with large squared errors.
The second describes a typical paired improvement. The vertical difference
between NMSE curves therefore need not equal the paired median. Reported
95\% intervals use 4,000 paired bootstrap resamples. They are pointwise
intervals, not simultaneous guarantees over all conditions. Averages over
SNR use equal weight for each stated operating point.

The main verification campaign contains 27,300 paired trials across
73 settings, including auxiliary diagnostics, and 2,400 bound evaluations
on the main-figure channels. The added baseline, split and position-error
studies contain 12,000 condition-trial evaluations from 4,800 independent
base realizations. Split fractions reuse the same observations, while
position-error strengths reuse the underlying paths, pilots and noise.
The separate update-budget check uses 300 further trials. Reference-strength
results reuse 2,800 trials from a separate campaign, and the rank-proxy
analysis replays EM-GS on 10,400 existing main-campaign trials. Replays and
paired repeats are not counted as new independent channels. Trial counts
and candidate grids are fixed before inspecting the corresponding results.

\subsection{Performance Against the Reference Curves}
\label{sec:main}

We first compare EM-GS and HS-GS with the CRLB and CCRB discussed in
Section~\ref{sec:crlb}. Fig.~\ref{fig:main}(a) shows the NMSE as the
aggregate SNR changes from $-5$ to $20$~dB. This experiment uses
$N=32$, $K=3$, $P=30$, and an RSR of $12$~dB.

HS-GS performs better than EM-GS at all six SNR values. The median
improvement across the paired trials ranges from $2.54$ to $2.97$~dB.
The average of these six median improvements is $2.70$~dB. The vertical
distance between the two NMSE curves ranges from $2.38$ to $2.90$~dB.
These values are slightly different because the paired median and the
NMSE curves summarize the trial errors in different ways.

For this experiment, the selected Hankel rank is smaller than the maximum
rank in all 2,400 trials. Therefore, the Hankel projection is active in
every trial.

At SNR values of $5$~dB and above, the EM-GS result differs from the
ordinary Rician CRLB by less than $0.108$~dB. This shows that the ordinary
CRLB is a useful reference for the EM-GS baseline.

The geometric CCRB is approximately $6.95$--$7.06$~dB below the ordinary
CRLB. It includes the local geometric multipath structure, while the
ordinary CRLB treats the channel coefficients as separate unknowns.
HS-GS improves on EM-GS but remains $4.37$--$4.85$~dB above the geometric
reference. These gaps describe the two reference curves and the measured
estimator performance under this model. They do not establish how much
further improvement a practical estimator can achieve.

The CRLB and CCRB are reference curves for unbiased estimators. EM-GS and
HS-GS are not assumed to be unbiased, so these curves are not treated as
absolute lower limits on their measured NMSE.

We define the vertical gain between the NMSE curves as

\begin{equation}
\Delta_{\mathrm{HS}}
=
\mathrm{NMSE}_{\mathrm{EM-GS,dB}}
-
\mathrm{NMSE}_{\mathrm{HS-GS,dB}}.
\end{equation}

A positive value of $\Delta_{\mathrm{HS}}$ means that HS-GS has a lower
NMSE than EM-GS.

\begin{figure*}[!t]
\centering
\includegraphics[width=\textwidth]{fig/figure2_performance_array.pdf}
\caption{Performance of EM-GS and HS-GS. (a) Ratio-of-sums NMSE versus
aggregate SNR for $N=32$, $P=30$ and RSR $=12$~dB, with the Rician CRLB
and geometric CCRB. Each SNR uses 400 paired trials. (b) HS-GS gain versus
array size for $P=10$ and $P=30$, averaged over six SNR values. Both the
ratio-of-sums and paired-median summaries are shown.}

\label{fig:main}
\end{figure*}

\subsection{Effect of Array Size}
\label{sec:aperture}

We next study how the number of receiver elements affects the benefit of
the Hankel constraint. We test

\begin{equation}
N\in\{8,16,32\}.
\end{equation}

For each array size, we use two pilot counts, $P=10$ and $P=30$, and six
SNR values. Each row in Table~\ref{tab:array} therefore summarizes
12 operating points.

\begin{table}[!t]
\caption{Effect of array size, averaged over 12 operating points. All
values except the array size are in dB.}
\label{tab:array}
\centering
\footnotesize
\setlength{\tabcolsep}{4pt}
\begin{tabular}{rrrr}
\toprule
$N$
& Ratio-of-sums gain
& Paired median
& EM-GS NMSE \\
\midrule
8  & $-0.095$ & $0.236$ & $-8.666$ \\
16 & $0.866$  & $1.044$ & $-8.786$ \\
32 & $2.938$  & $3.264$ & $-8.733$ \\
\bottomrule
\end{tabular}
\end{table}

Fig.~\ref{fig:main}(b) and Table~\ref{tab:array} show that the benefit of
HS-GS increases with the array size. The mean of the 12 per-setting
paired-median gains increases from
$0.236$~dB at $N=8$ to $3.264$~dB at $N=32$.

This behavior can be understood from the maximum possible Hankel rank,

\begin{equation}
r_{\max}
=
\left\lceil\frac{N}{2}\right\rceil.
\end{equation}

For the three tested array sizes,

\begin{equation}
\begin{array}{c|ccc}
N & 8 & 16 & 32 \\
\hline
r_{\max} & 4 & 8 & 16
\end{array}.
\end{equation}

When $N=8$, the Hankel matrix can have a maximum rank of only four. Since
the simulated channels normally contain between three and seven paths,
there may be little unused rank available. A low-rank projection can then
remove part of the actual channel instead of removing only noise and
estimation error.

When $N=32$, the maximum Hankel rank is 16. The same channel now occupies
a smaller part of the available rank. This gives HS-GS more low-rank
structure to use, which explains why its improvement is much larger.

The EM-GS baseline changes very little as the array size increases. Its
average NMSE changes by only $0.120$~dB between the tested array sizes.
Without the Hankel constraint, increasing the array size adds more
measurements, but it also adds more unknown channel coefficients.

At $N=8$, the mean of the 12 paired-median gains is $0.236$~dB, while
the mean of the 12 ratio-of-sums gains is $-0.095$~dB. These are averages
of operating-point summaries, not summaries obtained by pooling every
$N=8$ trial together. Within each operating point, large errors can have
more influence on the ratio of sums than on the paired median.

The Hankel projection is inactive in $42.94\%$ of the $N=8$ trials. In
these trials, HS-GS gives exactly the same estimate as EM-GS. An active
projection can still increase channel error in other trials. The fallback
rate therefore does not guarantee a positive gain, and the two aggregation
rules should be reported separately.

The projection is inactive in only $6.38\%$ of the $N=16$ trials and
$0.31\%$ of the $N=32$ trials. This provides further evidence that the
Hankel constraint becomes more useful as the array provides more structural
capacity.

\subsection{Effect of Channel Complexity}
\label{sec:pathcount}

We next study whether the number of propagation paths affects the benefit
of the Hankel constraint. We keep the array size fixed at $N=32$ and the
SNR at $5$~dB, and change only the number of paths $L$. The same value of
$L$ is used for all users.

The results show a clear trend. When the channel has only $L=2$ paths,
HS-GS improves over EM-GS by $7.377$~dB. As the number of paths increases,
the improvement becomes smaller. At $L=16$, the gain becomes
$-0.125$~dB. These are ratio-of-sums gains: at $L=16$, the measured
HS-GS NMSE is slightly higher than the EM-GS NMSE. Its paired-median
gain is zero. This result concerns the tested hard-rank procedure.

This can be understood from the rank of the Hankel matrix. Since $N=32$,
the maximum Hankel rank is

\begin{equation}
r_{\max}
=
\left\lceil
\frac{32}{2}
\right\rceil
=
16.
\end{equation}

Consider the two ends of the experiment. With only two paths, the channel
can have a Hankel rank of at most $2$, while the matrix can support a rank
of up to $16$. The channel is therefore strongly low rank, and there is
useful structure for HS-GS to use. In contrast, when $L=16$, the channel
can use most or all of the available Hankel rank. In this case, forcing a
low-rank approximation can remove an important part of the real channel
instead of only removing estimation error.

The EM-GS baseline changes by only $0.129$~dB over the same experiment.
Therefore, the large change in HS-GS gain mainly comes from the changing
usefulness of the Hankel constraint.

These results show that HS-GS works best when the channel
has relatively few important components compared with the maximum Hankel
rank. However, simply counting the number of paths is not always enough to
describe the real complexity of the channel. We take a closer look at this
next using effective rank.
\subsection{Effective Rank}
\label{sec:effrank}

Two channels can have the same number of paths but still have different
levels of complexity. For example, consider two channels that both have
four paths. In the first channel, all four paths may have similar strength.
In the second channel, two paths may be strong while the other two are very
weak. Although both channels have $L=4$, the second channel can behave
more like a channel with only two important components.

The singular values of the Hankel matrix help us measure this difference.
A large singular value represents an important component of the matrix,
while a very small singular value represents a weak component. For example,
consider

\begin{equation}
[10,\;9,\;8,\;7]
\end{equation}

and

\begin{equation}
[10,\;8,\;0.2,\;0.1].
\end{equation}

Both matrices can have rank $4$. However, in the second case, most of the
matrix can be explained by the first two components. Its real complexity
is therefore much smaller than its ordinary rank suggests.

To measure this, we use the effective rank proposed in
\cite{royvetterli2007},

\begin{equation}
r_{\mathrm{eff}}
=
\exp
\left(
-\sum_i p_i\ln p_i
\right),
\qquad
p_i
=
\frac{\sigma_i}{\sum_j\sigma_j},
\end{equation}

where $\sigma_i$ are the singular values of the Hankel matrix. The value
$p_i$ represents the relative importance of each singular value. If many
singular values are important, $r_{\mathrm{eff}}$ is large. If only a few
are important, $r_{\mathrm{eff}}$ is small.

We then compare the effective rank with the highest rank that the Hankel
matrix can support. We define

\begin{equation}
\rho
=
\frac{r_{\mathrm{eff}}}{r_{\max}}.
\end{equation}

We calculate this quantity separately for each user column. At an
operating point, the reported $\rho$ is the median over the user columns
and channel realizations. When resampling trials, we keep all users from
the same realization together. The values shown here are calculated from
the true noiseless channels and are used to explain the simulation results;
they are not inputs to rank selection.

The value $\rho$ can be understood as the fraction of the available Hankel
rank that is effectively being used by the channel. For example, if

\begin{equation}
r_{\mathrm{eff}}=4,
\qquad
r_{\max}=16,
\end{equation}

then

\begin{equation}
\rho
=
\frac{4}{16}
=
0.25.
\end{equation}

This means that the channel effectively uses about one quarter of the
available Hankel rank. We would expect the low-rank structure to be useful
in this case. If instead $r_{\mathrm{eff}}=14$, then

\begin{equation}
\rho
=
\frac{14}{16}
=
0.875.
\end{equation}

The channel now uses most of the available rank, so there is much less
low-rank structure left to use.

Fig.~\ref{fig:boundary} shows the HS-GS gain against $\rho$ for
$N\in\{16,32,64\}$. The same pilot count, RSR, iteration settings, SNR
distribution, and performance measure are used for all three array sizes.
The tested path counts are $L\in\{2,4,7,9,14\}$ for $N=16$,
$L\in\{2,4,\ldots,16\}$ for $N=32$, and
$L\in\{8,14,29,38,48\}$ for $N=64$. The plotted gain is the pooled
paired median over the stated random-SNR distribution. It is distinct
from the ratio-of-sums gain used for the NMSE curves.

The results show that the HS-GS gain becomes smaller as $\rho$ increases.
When $\rho$ is small, the channel uses only a small part of the available
Hankel rank and the hard-rank processing used by HS-GS provides a clear improvement. When
$\rho$ approaches approximately $0.5$, the improvement becomes small.

For example, at $N=32$ and $\rho=0.536$, the median gain is
$0.141$~dB, with a 95\% confidence interval of
$[0.090,0.211]$~dB. The gain is still positive, but it is small in
practice. Therefore, $\rho\approx0.5$ should be understood as a rough
description of the region where the tested HS-GS becomes less useful,
rather than as an exact threshold.

We also test what happens when a low-rank constraint is forced on a more
complex channel. For $N=16$ and $L=14$, the measured normalized effective
rank is

\begin{equation}
\rho=0.732.
\end{equation}

When rank seven is forced, the median HS-GS gain becomes
$-0.154$~dB, with a 95\% confidence interval of
$[-0.201,-0.111]$~dB. The negative gain shows that forcing a low-rank
constraint on a complex channel can make the estimate worse.

With adaptive rank selection, the pooled median gain is zero for the same
channels. This happens because the estimator can select the maximum rank
and return to EM-GS when a useful lower-rank constraint is not identified.
However, the rank-selection procedure does not guarantee an improvement in
every individual trial.

These results suggest that $\rho$ is a useful way to describe when the
tested HS-GS helps. However, $\rho$ is calculated here using the true
noiseless channel, which is not known at the receiver. We therefore use
$\rho$ to understand the simulation results rather than as a value used
directly by the estimator.

\begin{figure*}[!t]
\centering
\includegraphics[width=\textwidth]{fig/figure3_matched.pdf}
\caption{Pooled paired-median HS-GS gain against the normalized effective rank
$\rho=r_{\mathrm{eff}}/r_{\max}$ for $N\in\{16,32,64\}$. The same
simulation settings and performance measure are used for all three array
sizes. Error bars show pointwise 95\% intervals. The HS-GS gain becomes
small near half of the available rank in this tested design; this is not a
universal boundary for structural processing.}
\label{fig:boundary}
\end{figure*}

We also compare effective rank with a simpler description of complexity,
$L/r_{\max}$. In a separate prediction check, we fit a piecewise-linear
gain predictor to the eight matched $N=32$ operating points and evaluate
it on six points: $N=16$, $L\in\{2,4,7\}$ and $N=64$,
$L\in\{8,14,29\}$. The mean absolute prediction errors are
$2.067$~dB using $L$, $0.073$~dB using $L/r_{\max}$, and
$0.317$~dB using $\rho$. The 95\% bootstrap interval for the error of
the $\rho$ predictor minus that of the $L/r_{\max}$ predictor is
$[0.086,0.283]$~dB. Thus, normalized path count is the better predictor
in this check. This comparison is limited to the stated operating points
and interpolation rule, with some endpoint extrapolation. All six target
gains are positive, so it does not test the ability to detect harmful
structural processing. Both quantities also use channel information that
is unavailable at the receiver. We therefore treat $\rho$ as a useful
description of the results, without claiming that it is the best or a
receiver-ready predictor of gain.

\subsection{Additional Results}
\label{sec:kinv}

We test $(K,P)=(2,13),(3,20),(4,27)$ with $N=32$, five paths per user,
RSR $=10$~dB and the random SNR distribution in Table~\ref{tab:config}.
Among the trials with SNR at least $5$~dB, the paired-median HS-GS gains
are $3.381$, $3.509$ and $3.552$~dB. Their $0.170$~dB span describes
this joint change in users and pilots, with approximately fixed $P/(2K)$.
It does not isolate the effect of user count alone.

We also test two different path distributions at $N=32$ and $P=20$.
The clustered channel contains four clusters with ten rays each. Cluster
centres are uniform over $(-\pi/2,\pi/2)$, with ray offsets uniform within
$\pm5^\circ$ and angles clipped to the same range. The second channel has
ten independently drawn paths. The path gains are independent complex
Gaussian variables with variance equal to the reciprocal of the total
path count. In each trial, the resulting channel is rescaled to the
Frobenius norm of its paired baseline channel. RSR, random SNR and update
counts follow Table~\ref{tab:config}.

The paired-median gains are $1.182$~dB for the clustered channel, with
95\% interval $[0.995,1.398]$~dB, and $0.831$~dB for the ten-path
channel, with interval $[0.710,1.011]$~dB. These positive gains show that
HS-GS remains useful for these two path distributions. Both retain the
geometric ULA model, so this test alone does not establish robustness to
array or propagation-model errors.

\subsection{Projection Placement and Computation}
\label{sec:placement}

We compare Hankel processing after every EM-GS update with processing
only once after EM-GS has finished. The two variants use the same channel,
pilots, reference, noise and selected rank. There are 300 paired trials
at each of six SNRs, and four Cadzow sweeps per structural application.
The shared rank is selected using four-sweep interleaved candidates.
This isolates placement conditional on that rank; the variants are not
separately tuned. Equal sweeps per application also do not mean equal
total computation.

The mean of the six per-SNR paired-median interleaving gains is
$0.204$~dB, with a paired-bootstrap interval of $[0.169,0.242]$~dB.
Interleaving therefore gives a small positive improvement in this design.
This does not imply that it is always preferable under a fixed runtime
budget or that the two estimates are otherwise identical.

A separate placement timing study uses $P=20$, SNR $=5$~dB,
RSR $=10$~dB, 100 final updates, 25 candidate updates and four sweeps.
It measures three independent trials per array size on an Apple M2
with eight physical cores and 8~GiB RAM, using single-thread settings.
At $N=32$, median total times are $3.548$~s for interleaving and
$3.174$~s for final-only processing. Rank selection accounts for
$79.3\%$ of the interleaved total at $N=32$ and $88.5\%$ at $N=64$.
These measurements describe that implementation and machine.

\subsection{Comparison with Other Structural Processing}
\label{sec:revision-structured-comparisons}

We compare six classical estimators: EM-GS, HS-GS, unstructured Rician
gradient descent, Rician gradient descent with Hankel processing, EM-GS
with smooth singular-value filtering, and the one-dimensional linearized
least-squares objective of Xu et al.~\cite{xu2025}. All methods receive
the same exact magnitude observations. The linearized method is a
comparison estimator; it does not change how the observations are generated.
Both nonlinear data updates use the Rician likelihood. Omitting terms
independent of $\mathbf G$, its negative log-likelihood is
\begin{equation}
\mathcal L(\mathbf G)=\sum_{n,p}\left[
\frac{a_{np}^{2}}{\sigma^2}
-\log I_0\!\left(\frac{2Z_{np}a_{np}}{\sigma^2}\right)\right],
\end{equation}
where $a_{np}=|(\mathbf G\mathbf S+\mathbf B)_{np}|$.
Scaled Bessel functions are used in the numerical calculation. Structural
correction is not assumed to decrease this likelihood at every update.
The Rician gradient method uses
\begin{equation}
\mathbf G^{(t+1)}_{\rm dc}
=\mathbf G^{(t)}-
\frac{(\mathbf M^{(t)}-\mathbf Y^{(t)})\mathbf S^H}
     {\|\mathbf S\|_2^2},
\end{equation}
where $\mathbf Y^{(t)}$ is the Rician conditional field in
Section~\ref{sec:revised-estimator}. The structured version follows this
update with the same per-user Cadzow processing as HS-GS. It is an adapted
Rician--Hankel control, rather than Xu et al.'s two-dimensional projected
gradient method.

For the smooth-filtering control, each singular value $s_i$ is replaced by
\begin{equation}
\widehat s_i=\frac{s_i^3}{s_i^2+\tau^2},
\qquad
\tau=f\sqrt{\frac{1}{q}\sum_{i=1}^{q}s_i^2},
\end{equation}
where $q=r_{\max}$ and
$f\in\{0,0.05,0.1,0.2,0.4,0.8,1.6\}$.
The scale is recomputed from the current per-user Hankel spectrum in each
of the four sweeps. A zero factor bypasses filtering. This is an empirical
smooth filter, not OptShrink or a nuclear-norm solver. Each structured
method selects its own rank or factor using the same held-out squared
magnitude residual. The selected parameter is shared across users, and
the final estimate is reinitialized and fitted using all pilots.

We use $N=32$, $K=3$, $P=30$, RSR $=12$~dB,
$L\in\{3,7,14\}$ paths per user and SNRs of 0 and 10~dB. Each of the
six settings uses 400 paired trials. The iterative methods use 100 final
updates and 25 updates per candidate, with four sweeps per structural
application. These are fixed update budgets, not convergence guarantees.
The gains below compare ratio-of-sums NMSE; intervals are pointwise 95\%
paired-bootstrap intervals from 4,000 resamples.

Fig.~\ref{fig:revision-baselines} shows that HS-GS is not the best method
in every setting. Smooth filtering improves on HS-GS by
$0.623$--$1.413$~dB in five settings, with the interval above zero in each.
The remaining case, $L=3$ and SNR $=10$~dB, gives a gain of
$0.005$~dB with interval $[-0.180,0.191]$~dB.
Rician--Hankel processing improves on HS-GS by $0.171$~dB
$[0.078,0.274]$ at $L=14$ and SNR $=0$~dB, but is worse by
$0.284$ and $0.222$~dB at SNR $=10$~dB for $L=3$ and $L=7$,
respectively. The two unstructured nonlinear methods have similar NMSE
in these tests. The linearized least-squares method is less accurate than
EM-GS at all six settings under this reference strength.

An important distinction appears at $L=14$ and SNR $=0$~dB. HS-GS has
a gain of $-0.147$~dB $[-0.251,-0.046]$ over EM-GS, while smooth
filtering has a gain of $1.266$~dB $[1.152,1.380]$. Thus, a loss for
HS-GS does not show that every use of Hankel structure is unhelpful. The
observed benefit also depends on how the structure is applied and how its
parameter is selected.

We check the final update budget using 300 additional paired trials,
with 50 trials for each combination of $L\in\{3,7,14\}$ and
SNR $\in\{0,10\}$~dB. The remaining settings are unchanged. For each
method, the rank or shrinkage parameter is selected once using 25
candidate updates, and the same final trajectory is recorded at 100 and
300 updates. Across the five iterative methods and six operating points,
the absolute change in ratio-of-sums NMSE is less than $0.06$~dB.
This check supports the chosen final update budget in the tested
conditions, but it does not establish equivalence, global convergence, or
the best possible rank-selection budget. Some estimates are still changing
slightly after 300 updates. Parameters are not selected again using the
longer trajectories or the true channel errors.

\begin{figure*}[!t]
\centering
\includegraphics[width=\textwidth]{fig/revision_structured_baselines.pdf}
\caption{Six classical estimators on the same exact magnitude observations.
Each point uses 400 paired trials at $N=32$, $K=3$, $P=30$ and
RSR $=12$~dB. The two panels use SNRs of 0 and 10~dB.
Error bars are pointwise 95\% paired-bootstrap intervals.}
\label{fig:revision-baselines}
\end{figure*}

\subsection{Sensitivity to Array Position Errors}
\label{sec:revision-position-mismatch}

To test one departure from the ideal ULA model, we perturb the element
positions according to
\begin{equation}
\frac{x_n}{\lambda}=\frac{n-1}{2}+s_x u_n,
\qquad u_n\sim\mathcal N(0,1),
\end{equation}
with independent $u_n$ and $s_x\in\{0,0.005,0.01,0.02\}$.
These are declared stress levels,
not measured hardware tolerances. The estimator continues to apply its
ULA Hankel processing, and channel error is measured against the actual
perturbed channel. The reference field remains known and calibrated.

We use $N=32$, $K=3$, $P=30$, SNR $=5$~dB, RSR $=12$~dB and
$L\in\{3,12\}$ paths per user. Each setting uses 400 paired trials and
the same 100/25 update counts, four sweeps and 30\% validation fraction.
For a given trial, the path parameters, pilots, thermal-noise draw and
normalized offsets $u_n$ are shared across the four strengths. We do not
renormalize the realized channel after perturbing the positions; SNR
retains its nominal statistical-power definition.

As shown in Fig.~\ref{fig:revision-position}, the HS-GS gain decreases
as the position errors increase. For $L=3$, its ratio-of-sums gain over
EM-GS changes from $4.955$~dB without mismatch to $3.917$~dB at
$s_x=0.02$, with interval $[3.748,4.075]$~dB at the latter point.
For $L=12$, the gain decreases from $0.264$ to $0.050$~dB; the final
interval, $[-0.029,0.125]$~dB, includes zero. Thus, the tested sparse
case retains a benefit, while the smaller benefit in the more complex
case becomes uncertain. This experiment does not establish robustness to
unknown reference calibration, mutual coupling, near-field propagation
or correlated noise.

\begin{figure*}[!t]
\centering
\includegraphics[width=\textwidth]{fig/revision_position_mismatch.pdf}
\caption{Position-error stress test with a known calibrated reference and
the actual perturbed channel as the estimation target. Each point uses
400 trials; trials are paired across mismatch strengths within each path
count. Error bars are pointwise 95\% paired-bootstrap intervals from
4,000 resamples.}
\label{fig:revision-position}
\end{figure*}

\subsection{Runtime of the Comparison Methods}
\label{sec:comparison-runtime}

We also time the six classical methods in a separate serial experiment. This uses a vectorized implementation on the same machine,
$N=32$, $K=3$, $P=30$, $L=7$, SNR $=5$~dB and RSR $=12$~dB.
Table~\ref{tab:revision-serial} reports medians over 20 fresh worlds with
single-thread CPU execution and double precision. The iterative methods
use 100 final updates, 25 updates per candidate and four structural
sweeps. Timings include initialization, parameter selection where needed,
and the final fit, but exclude channel and observation generation.

\begin{table}[!t]
\caption{Serial runtime of the classical comparator implementations.}
\label{tab:revision-serial}
\centering
\footnotesize
\begin{tabular}{lr}
\toprule
Method & Median total time (ms)\\
\midrule
EM-GS & 9.943\\
HS-GS & 348.536\\
Rician GD & 10.015\\
Rician GD + Hankel & 346.500\\
EM-GS + smooth filtering & 200.799\\
Xu 1D linearized LS & 0.647\\
\bottomrule
\end{tabular}
\end{table}

The methods are timed in a fixed order for each world: HS-GS,
Rician--Hankel, smooth filtering, EM-GS, Rician GD and linearized LS.
No warm-up worlds are discarded. The median selection share is 76.0\%
for HS-GS and 56.4\% for smooth filtering; the latter searches seven
threshold factors rather than 16 ranks. These are measured costs for one
implementation, machine and operating point.
They are separate from the placement timings in Section~\ref{sec:placement} at
$P=20$ and RSR $=10$~dB, so a difference between the two timing studies
must not be interpreted as an algorithmic speedup.

\subsection{Reference Strength, Pilot Splitting, and Rank Proxies}
\label{sec:sensitivity-summary}

Further tests address practical choices behind HS-GS. At $N=32$, $P=30$
and SNR $=5$~dB, the ratio-of-sums gain remains positive across the seven
RSR values from 0 to 25~dB, ranging from $1.83$ to $2.51$~dB. A weak
reference still worsens the absolute NMSE of both methods. Thus, a positive
structural gain does not remove the difficulty caused by a weak reference.
Appendix~\ref{app:additional-sensitivity} gives the full reference sweep.

Validation fractions of 10\%, 20\%, 30\% and 40\% produce different ranks
and errors. For example, at $P=30$ and SNR $=10$~dB, the 10\% split is
$0.602$~dB worse than the 30\% split. At $P=10$ and the same SNR, the
40\% split is $0.461$~dB worse than 30\%. All pilots are reused for the
final estimate. The best fraction changes with the setting, so 30\% is a
declared choice, not a generally optimal one.

An effective-rank proxy computed from the unstructured EM-GS estimate
tracks the true effective rank more closely at higher SNR. However, noise
raises its value. At $N=32$ and SNR $=5$~dB, the mean proxy is $0.577$
while the corresponding true-channel mean is $0.244$, even though the
paired-median HS-GS gain is $2.544$~dB. The proxy study in
Appendix~\ref{app:additional-sensitivity} is descriptive; it does not
validate a per-trial threshold or replace held-out rank selection.

\section{Limitations and Conclusion}

This work studies how spatial Hankel processing affects magnitude-only
channel estimation for the stated Rydberg receiver model. In the main
experiment, HS-GS improves on EM-GS by an average of $2.70$~dB across the
six per-SNR paired medians. Its benefit increases with array capacity and
generally decreases as channel complexity increases. Effective rank helps
explain this trend, but it is not a universal predictor or a validated
receiver-side decision rule.

The stronger comparisons change how this trend should be interpreted.
Smooth singular-value filtering improves on hard-rank HS-GS in five of
six new comparison settings, by $0.623$--$1.413$~dB with positive
pointwise intervals. It can improve estimation even where HS-GS loses
accuracy. Therefore, the point where the tested HS-GS becomes less useful
is a property of its processing and rank-selection procedure, not a
universal limit of Hankel structure. Hard truncation, the validation split
and the update budget should all be treated as design choices.

Interleaving provides a small improvement over final-only processing under
the shared-rank comparison, while rank selection remains the main measured
computational cost. Reference strength and array-position errors also affect
the benefit. The practical conclusion is to evaluate the structural rule
and its parameter selection together for the intended channel and pilot
conditions, including cases where the rule can cause harm.
\textit{Limitations.} All results are from simulations; we have not tested
measured Rydberg receiver data. The additional path distributions retain
the geometric channel model. The position-error study tests one departure
from an ideal ULA with a known calibrated reference, but does not cover
unknown calibration, mutual coupling, near-field propagation or correlated
noise. Its offset levels are stress-test settings rather than hardware
tolerances.

The true effective rank remains an analysis quantity that is unavailable
at the receiver. A proxy computed from EM-GS provides useful descriptive
information, especially at higher SNR, but noise changes its value.
Neither the true-rank plots nor the proxy study establish a universal
cutoff or a validated per-trial switching rule. The comparison with smooth
filtering also shows that the useful region depends on the structural
operation and its parameter selection.

The comparator results use the stated parameter grids and update budgets;
they do not establish optimality over all structured estimators. The longer
update check is limited to its tested conditions, and the timing results
describe one implementation and machine. Finally, the CRLB and CCRB are
local unbiased-estimation references. The tested estimators need not be
unbiased, so the bounds are not absolute lower limits on their NMSE.

\appendices
\section{Real-Coordinate Fisher Information and Geometric CCRB}
\label{app:revision-ccrb}

\subsection{Channel coordinates and measurement information}

For receive element $n$, define
\begin{equation}
\mathbf x_n=
 [\Re G_{n1},\ldots,\Re G_{nK},
  \Im G_{n1},\ldots,\Im G_{nK}]^T.
\end{equation}
The full real vector is
$\mathbf x=[\mathbf x_1^T,\ldots,\mathbf x_N^T]^T\in\mathbb R^{2NK}$.
This ordering gives
$\|\widehat{\mathbf x}-\mathbf x\|_2^2
=\|\widehat{\mathbf G}-\mathbf G\|_F^2$.
We condition on the known $\mathbf S$, $\mathbf B$ and $\sigma^2$.
The noise convention is $\mathbb E|W_{np}|^2=\sigma^2$, so each real
component has variance $\sigma^2/2$.

Write $\mu_{np}=(\mathbf G\mathbf S+\mathbf B)_{np}$ and
$a_{np}=|\mu_{np}|$. For $z\geq0$, the amplitude density is
\begin{equation}
p(z\mid a)=\frac{2z}{\sigma^2}
 \exp\!\left(-\frac{z^2+a^2}{\sigma^2}\right)
 I_0\!\left(\frac{2za}{\sigma^2}\right).
\end{equation}
With $R(u)=I_1(u)/I_0(u)$, its amplitude score is
\begin{equation}
\frac{\partial\log p(z\mid a)}{\partial a}
=\frac{2}{\sigma^2}
 \left[zR\!\left(\frac{2za}{\sigma^2}\right)-a\right].
\end{equation}
Using $\mathbb E[zR(2za/\sigma^2)]=a$, the scalar information is
\begin{align}
\mathcal I_a(a)
&=\frac{4}{\sigma^4}
 \left[\mathbb E\!\left\{z^2R^2\!\left(\frac{2za}{\sigma^2}\right)\right\}
 -a^2\right],
\label{eq:revision-scalar-information}\\
\beta(a)&=\mathcal I_a(a)/4.
\end{align}

For $\mu_{np}\ne0$, put
\begin{equation}
\mathbf c_{np}=e^{-\jmath\arg\mu_{np}}\mathbf S_{:,p},
\qquad
\mathbf v_{np}=
\begin{bmatrix}\Re\mathbf c_{np}\\-\Im\mathbf c_{np}\end{bmatrix}.
\end{equation}
Here $\mathbf v_{np}=\partial a_{np}/\partial\mathbf x_n$.
Independence of the measurement noises gives
\begin{align}
\mathbf J_n&=\sum_{p=1}^{P}
  \mathcal I_a(a_{np})\mathbf v_{np}\mathbf v_{np}^T,\\
\mathbf J&=\operatorname{blkdiag}
 (\mathbf J_1,\ldots,\mathbf J_N).
\label{eq:revision-real-fisher}
\end{align}
Each nonzero measurement term has rank one in real coordinates. At
$\mu_{np}=0$, the contribution is defined by its continuous zero limit;
an arbitrary choice of phase supplies no information.

\subsection{Geometric tangent space}

For each user $k$ and path $\ell$, take the real parameters
$(\psi_{\ell,k},\Re\alpha_{\ell,k},\Im\alpha_{\ell,k})$.
Collect them in $\boldsymbol\eta$ and set
$\mathbf D=\partial\mathbf x/\partial\boldsymbol\eta$.
Writing $v_n(\psi)=e^{-\jmath(n-1)\psi}$, the complex derivatives for
one path are
\begin{align}
\frac{\partial g_k[n]}{\partial\psi_{\ell,k}}
 &=-\jmath(n-1)\alpha_{\ell,k}v_n(\psi_{\ell,k}),\\
\frac{\partial g_k[n]}{\partial\Re\alpha_{\ell,k}}
 &=v_n(\psi_{\ell,k}),\\
\frac{\partial g_k[n]}{\partial\Im\alpha_{\ell,k}}
 &=\jmath v_n(\psi_{\ell,k}).
\end{align}
Their real and imaginary parts fill the corresponding entries of
$\mathbf x_n$; entries belonging to other users are zero. Using spatial
frequency $\psi$ follows the channel generator directly and avoids a
coordinate singularity from $\psi=\pi\sin\theta$ at endfire.

Take a thin SVD $\mathbf D=\mathbf U_D\boldsymbol\Sigma_D\mathbf V_D^T$.
The implementation retains the columns of $\mathbf U_D$ whose singular
values exceed $10^{-9}$ times the largest singular value. These columns
form $\mathbf U$, an orthonormal basis for the numerically resolved range
of $\mathbf D$. This construction avoids inverting a redundant physical
parameterization when $\mathbf D$ is not full column rank.

For locally unbiased estimation along this tangent space,
\begin{equation}
\mathbf C_{\rm geom}
=\mathbf U(\mathbf U^T\mathbf J\mathbf U)^{-1}\mathbf U^T.
\end{equation}
The gains and spatial frequencies are varied jointly in $\mathbf D$;
neither is held fixed as known side information. This calculation includes
no unknown reference-field, calibration or noise-variance parameters.
Adding such nuisance parameters would require the corresponding enlarged
information matrix. The true path counts and parameter point specify the
local reference, rather than inputs available to HS-GS.

The displayed inverses require positive definite information on the
relevant space. A singular restricted matrix indicates an unresolved
local direction and is not replaced by a regularized finite bound.
The known reference can remove phase ambiguity, but its presence alone
does not guarantee identifiability for every pilot design. Coincident or
zero-gain paths can also invalidate a regular local parameterization.
The stated numerical rank tolerance must therefore accompany the results.

\subsection{Numerical evaluation and normalization}

For trial $i$, the two scalar reference errors are
\begin{align}
b_{{\rm unc},i}&=\sum_{n=1}^{N}\operatorname{tr}(\mathbf J_{n,i}^{-1}),\\
b_{{\rm geom},i}&=
 \operatorname{tr}\!\left[(\mathbf U_i^T\mathbf J_i\mathbf U_i)^{-1}\right].
\end{align}
The second identity uses $\mathbf U_i^T\mathbf U_i=\mathbf I$.
For either reference, the plotted value is
\begin{equation}
10\log_{10}
 \frac{\sum_i b_i}{\sum_i\|\mathbf G_i\|_F^2}.
\label{eq:revision-bound-normalization}
\end{equation}
The same channel realizations and energy denominators are used for the
estimator curves. This is a ratio of sums, not an average of trial-wise
dB values.

The scalar expectation is evaluated with scaled Bessel functions and
quadrature over $[\max(0,a-8\sigma),a+8\sigma]$, where
$\sigma=\sqrt{\sigma^2}$. The quadrature uses absolute tolerance
$10^{-14}$ and relative tolerance $10^{-12}$, splitting at $a$ when it
lies inside the integration window. At each main-figure SNR, 800 equally
spaced amplitude nodes cover $0.99a_{\min}$ to $1.01a_{\max}$ over that
SNR's worlds. Piecewise cubic Hermite interpolation supplies the scalar
weights at the measurement amplitudes. The implementation checks 200
sampled amplitudes against direct quadrature and requires maximum relative
error below $10^{-4}$. The tangent derivatives are checked by finite
differences. It also checks that the constrained trace is positive and
no larger than the unconstrained trace, and records the condition number
of the restricted information matrix. No diagonal ridge is added.

\section{Additional Sensitivity Results}
\label{app:additional-sensitivity}
\subsection{Effect of Reference Strength}
\label{sec:revision-rsr}

We also test how the reference strength affects the benefit of the Hankel
constraint. We use $N=32$, $K=3$, $P=30$, SNR $=5$~dB, and independently
sampled path counts $L_k\in\{3,\ldots,7\}$. The single-user RSR takes
the values $0$, $5$, $10$, $12$, $15$, $20$, and $25$~dB. Both estimators
use 100 final updates. HS-GS uses 25 updates per candidate rank, four
Cadzow sweeps per application, and the same 30\% validation split followed
by a final fit using all pilots.

This supplementary analysis reuses 2,800 stored paired trials from a
separate simulation campaign, with 400 trials at each RSR. These are
distinct from the 27,300 main-campaign trials reported above. No new Monte
Carlo realizations were generated for this analysis. We use only the
classical EM-GS and HS-GS results. Pointwise 95\% confidence intervals
use 4,000 paired bootstrap resamples.

Fig.~\ref{fig:revision-rsr} shows that HS-GS improves on EM-GS at all
seven RSR values in this experiment. The ratio-of-sums gain ranges from
$1.83$ to $2.51$~dB, and each corresponding confidence interval is above
zero. At RSR $=0$~dB, the gain is $1.83$~dB, with a 95\% interval of
$[1.40,2.26]$~dB. The paired-median gain at this point is $2.83$~dB;
these are different summaries of the same paired errors.

The weak reference still makes channel estimation harder for both
methods. At RSR $=0$~dB, the EM-GS and HS-GS NMSEs are $-6.19$ and
$-8.01$~dB. At RSR $=12$~dB, they are $-10.73$ and $-12.93$~dB.
Thus, the structural gain remains useful at the weakest tested reference,
but it does not remove the loss caused by that weak reference. The gain
does not increase steadily with RSR. These findings apply to the tested
array size, SNR, and channel distribution.

\begin{figure*}[!t]
\centering
\includegraphics[width=\textwidth]{fig/rsr_classical.pdf}
\caption{Effect of single-user RSR for $N=32$, $K=3$, $P=30$, and
SNR $=5$~dB. Left: ratio-of-sums NMSE. Right: HS-GS gain measured by the
ratio of summed errors and by the paired median. Shaded bands and error
bars are pointwise 95\% paired-bootstrap intervals. Each RSR uses 400
paired trials from a separate stored campaign.}
\label{fig:revision-rsr}
\end{figure*}

\subsection{Sensitivity to the Validation Fraction}
\label{sec:revision-validation-fraction}

We test validation fractions of 10\%, 20\%, 30\% and 40\% at
$P\in\{10,30\}$ and SNRs of 0 and 10~dB. The other settings are
$N=32$, $K=3$, RSR $=12$~dB and independently drawn three-to-seven
paths per user. We retain the 100 final updates, 25 updates per candidate
and four Cadzow sweeps. For each pilot-count/SNR pair, the four fractions
use the same 400 channel, pilot and noise realizations. The validation
sets are nested trailing pilot columns. All $P$ pilots are reused for
every final fit, so the comparison changes the rank-selection split
rather than the total transmitted pilot count.

Fig.~\ref{fig:revision-holdout} shows that the split matters. For example,
at $P=30$ and SNR $=10$~dB, using 10\% instead of 30\% gives a paired
ratio-of-sums gain of $-0.602$~dB, with interval
$[-0.832,-0.383]$~dB. At $P=10$ and SNR $=10$~dB, increasing the
fraction from 30\% to 40\% gives $-0.461$~dB
$[-0.744,-0.176]$. Negative values mean that the changed split is worse
than the 30\% reference. Neither a very small nor a larger validation
set is therefore automatically better.

The selected ranks also change. At $P=10$ and SNR $=0$~dB, the median
selected rank falls from three with a 10\% split to one with a 40\%
split. Across all 16 settings, full-rank fallback occurs in 0--1.5\% of
the trials. HS-GS still improves on EM-GS in each setting, but the results
do not identify one validation fraction that is best in all conditions.

\begin{figure*}[!t]
\centering
\includegraphics[width=0.80\textwidth]{fig/revision_heldout_fraction.pdf}
\caption{HS-GS ratio-of-sums gain over EM-GS as the validation fraction
changes. The fractions share 400 paired worlds within each panel, and
all pilots are used for the final estimate. Error bars are pointwise
95\% paired-bootstrap intervals from 4,000 resamples.}
\label{fig:revision-holdout}
\end{figure*}

\begin{figure*}[!t]
\centering
\includegraphics[width=\textwidth]{fig/rho_proxy_diagnostics.pdf}
\caption{Effective rank estimated from unstructured EM-GS.
(a) Correlation between the observable proxy and true effective rank,
using trial-level user means at each fixed SNR.
(b) Correlation with paired HS-GS gain across $L\in\{2,4,\ldots,16\}$
within fixed 5~dB SNR bands. The two true-channel quantities are
evaluation-only references. Error bars are pointwise 95\% intervals
from 4,000 paired trial-bootstrap resamples.}
\label{fig:revision-observable-rank}
\end{figure*}

\subsection{Effective Rank Estimated from Measurements}
\label{sec:revision-observable-rank}

The true channel is unavailable at the receiver, so the effective rank
used above cannot directly control the estimator. We therefore test a
simple observable proxy: calculate the effective rank from the final
unstructured EM-GS estimate. For user $k$ in trial $i$, define
\begin{equation}
\begin{aligned}
\widehat\rho_{{\rm EM},k,i}
&=\frac{r_{\rm eff}\!\left(\mathcal H(
  \widehat{\mathbf g}_{{\rm EM},k,i})\right)}{r_{\max}},\\
\overline{\widehat\rho}_{{\rm EM},i}
&=\frac{1}{K}\sum_{k=1}^{K}\widehat\rho_{{\rm EM},k,i}.
\end{aligned}
\label{eq:revision-observable-rho}
\end{equation}
We compare this trial-level mean with the corresponding mean from the
true user channels and with the paired HS-GS gain. We use EM-GS estimates
so that the proxy is measured before a Hankel correction changes the
singular values. No true channel or path count is used to form the proxy.

This analysis reuses 10,400 existing main-campaign trials: 7,200 trials from the
array-size study with $P=30$, and 3,200 trials from the matched $N=32$
complexity study with $P=20$. We repeat only EM-GS using the original
channels, observations, seeds, and update counts. Every reproduced EM-GS
squared error agrees exactly with its stored value in this implementation.
These replays are not additional independent Monte Carlo trials.

Fig.~\ref{fig:revision-observable-rank}(a) shows that the proxy tracks
true effective rank more closely as SNR increases. For $N=32$, the
Spearman correlation rises from $0.202$ at $-5$~dB, with a 95\% interval
of $[0.100,0.297]$, to $0.914$ at $20$~dB, with an interval of
$[0.890,0.931]$. The intervals use 4,000 bootstrap resamples of entire
trials, keeping their user averages and paired errors together.

However, estimation noise raises the proxy substantially. For $N=32$
and SNR $=5$~dB, its mean across trials is $0.577$, whereas the
corresponding true-channel mean is $0.244$. HS-GS still gives a positive
paired-median gain of $2.544$~dB. These are different summaries: the first
two numbers are means of effective rank, while the last is a median of
paired estimation gains. The upward shift warns against applying a cutoff
near $0.5$ to this proxy without calibration. This comparison of averages
does not evaluate a per-trial cutoff.

We also examine the relation between the proxy and HS-GS gain across
path counts. We use the matched $N=32$, $P=20$, RSR $=10$~dB study and
group its random SNRs into fixed 5~dB bands. In the $[-10,-5)$~dB band,
the proxy--gain correlation is $-0.098$, with a 95\% interval of
$[-0.183,-0.012]$. In the $[15,20]$~dB band, it is $-0.797$, with an
interval of $[-0.836,-0.752]$. Negative correlation means that larger
estimated effective rank tends to accompany a smaller structural gain.
The true effective rank and normalized true path count are also shown
as diagnostic references; neither is available to the receiver.

The proxy therefore provides some practical information, especially at
higher SNR. It is not a calibrated replacement for the true effective
rank, and no per-trial rank-selection rule or switching threshold is
validated here. The SNR bands reduce, but do not remove, the influence
of SNR. In this diagnostic we use means over users within each trial;
the effective-rank curves in Section~\ref{sec:effrank} instead summarize their operating
points by medians over user columns and trials.



\clearpage
\IEEEtriggeratref{9}
\begin{thebibliography}{99}

\bibitem{cui2025}
M.~Cui, Q.~Zeng, and K.~Huang,
``Towards atomic MIMO receivers,''
\emph{IEEE J. Sel. Areas Commun.},
vol.~43, no.~3, pp.~659--673, 2025.

\bibitem{gong2025mag}
T.~Gong \emph{et al.},
``Rydberg atomic quantum receivers for classical wireless communication and sensing,''
\emph{IEEE Wireless Commun.},
vol.~32, no.~5, pp.~90--100, 2025.

\bibitem{gerchberg1972}
R.~W.~Gerchberg and W.~O.~Saxton,
``A practical algorithm for the determination of phase from image and diffraction plane pictures,''
\emph{Optik},
vol.~35, pp.~237--246, 1972.

\bibitem{xu2025}
B.~Xu, J.~Zhang, Z.~Chen, B.~Cheng, Z.~Liu, Y.-C.~Wu, and B.~Ai,
``Channel estimation for Rydberg atomic receivers,''
\emph{IEEE Wireless Commun. Lett.},
vol.~14, no.~9, pp.~2957--2961, 2025.

\bibitem{xiao2026}
J.~Xiao, J.~Wang, M.~Zeng, H.~Xu, X.~Li, and A.~Nallanathan,
``Channel estimation for Rydberg atomic quantum receivers:
unrolled phase retrieval from holographic snapshots,''
\emph{IEEE Signal Process. Lett.},
vol.~33, pp.~1696--1700, 2026.

\bibitem{cadzow1988}
J.~A.~Cadzow,
``Signal enhancement---a composite property mapping algorithm,''
\emph{IEEE Trans. Acoust., Speech, Signal Process.},
vol.~36, no.~1, pp.~49--62, 1988.

\bibitem{markovsky2008}
I.~Markovsky,
``Structured low-rank approximation and its applications,''
\emph{Automatica},
vol.~44, no.~4, pp.~891--909, 2008.

\bibitem{chen2014emac}
Y.~Chen and Y.~Chi,
``Robust spectral compressed sensing via structured matrix completion,''
\emph{IEEE Trans. Inf. Theory},
vol.~60, no.~10, pp.~6576--6601, 2014.

\bibitem{jin2016aloha}
K.~H.~Jin, D.~Lee, and J.~C.~Ye,
``A general framework for compressed sensing and parallel MRI using annihilating filter based low-rank Hankel matrix,''
\emph{IEEE Trans. Comput. Imag.},
vol.~2, no.~4, pp.~480--495, 2016.

\bibitem{haldar2014loraks}
J.~P.~Haldar,
``Low-rank modeling of local $k$-space neighborhoods (LORAKS) for constrained MRI,''
\emph{IEEE Trans. Med. Imag.},
vol.~33, no.~3, pp.~668--681, 2014.

\bibitem{nadakuditi2014optshrink}
R.~R.~Nadakuditi,
``OptShrink: An algorithm for improved low-rank signal matrix denoising by
optimal, data-driven singular value shrinkage,''
\emph{IEEE Trans. Inf. Theory}, vol.~60, no.~5, pp.~3002--3018, 2014,
doi: 10.1109/TIT.2014.2311661.

\bibitem{balan2015}
R.~Balan,
``The Fisher information matrix and the Cram\'er--Rao bound in a non-AWGN model for the phase retrieval problem,''
in \emph{Proc. Int. Conf. Sampling Theory Appl. (SampTA)},
2015.

\bibitem{gorman1990}
J.~D.~Gorman and A.~O.~Hero,
``Lower bounds for parametric estimation with constraints,''
\emph{IEEE Trans. Inf. Theory},
vol.~36, no.~6, pp.~1285--1301, 1990.

\bibitem{stoica1998}
P.~Stoica and B.~C.~Ng,
``On the Cram\'er--Rao bound under parametric constraints,''
\emph{IEEE Signal Process. Lett.},
vol.~5, no.~7, pp.~177--179, 1998.

\bibitem{royvetterli2007}
O.~Roy and M.~Vetterli,
``The effective rank: a measure of effective dimensionality,''
in \emph{Proc. Eur. Signal Process. Conf. (EUSIPCO)},
2007, pp.~606--610.

\end{thebibliography}

\end{document}
